\chapter{Complex Numbers}

Complex numbers are an extension of the real numbers, allowing us to work with quantities that have both a real and an imaginary part. A complex number is typically expressed in the form \( z = a + bi \), where:
\begin{itemize}
    \item \( a =\Re(z) \) is the real part of the complex number.
    \item \( b =\Im(z) \) is the imaginary part of the complex number.
    \item \( i \) is the imaginary unit, defined by the property \( i^2 = -1 \).
\end{itemize}

They can be expressed using diferent notations:
\begin{description}
    \item[Cartesian/Algebraic Form:] \( z = a + bi \)
    \item[Polar/Trigonometric Form:] \( z = r(\cos \theta + i \sin \theta) \)
    
    The idea behind the polar form is to use polar coordinates to represent the point $(a, b)$ in the complex plane. Here, \( r \) is the magnitude (or modulus) of the complex number, and \( \theta \) is the argument (or angle) of the complex number, measured from the positive real axis.
    \item[Exponential Form:] \( z = re^{i\theta} \)
    
    Using the Euler's formula~\eqref{eq:Euler} we can express the complex number in exponential form, which is often more convenient for multiplication and division of complex numbers.
    \begin{equation}
        e^{i\theta} = \cos \theta + i \sin \theta
        \label{eq:Euler}
    \end{equation}

    Using the exponential form, the Moive's theorem~\eqref{eq:Moivre} can be easily derived, which states that for any complex number \( z = re^{i\theta} \) and any integer \( n \):
    \begin{equation} 
        z^n = r^n \left( \cos(n\theta) + i \sin(n\theta) \right) = r^n e^{in\theta}
        \label{eq:Moivre}
    \end{equation}
\end{description}


The reader should be familiar with the basic operations on complex numbers and with the different representations of complex numbers, as well as with the properties of the complex plane.